\documentclass[11pt]{article}
\usepackage{amsmath,amssymb,amsthm}
\usepackage[margin=1in]{geometry}
\usepackage[colorlinks=true,linkcolor=blue,citecolor=blue,urlcolor=blue]{hyperref}

\newtheorem{theorem}{Theorem}
\newtheorem{proposition}{Proposition}
\newtheorem{lemma}{Lemma}
\newtheorem{corollary}{Corollary}
\newtheorem{remark}{Remark}
\theoremstyle{definition}

\DeclareMathOperator{\Nm}{N}
\newcommand{\Z}{\mathbb{Z}}
\newcommand{\Q}{\mathbb{Q}}
\newcommand{\R}{\mathbb{R}}
\newcommand{\norm}[1]{\lVert #1 \rVert}

\title{Powerful numbers between consecutive squares}
\author{Scott D. Hughes}
\date{\today}

\begin{document}
\maketitle

\begin{abstract}
Let $h(n)$ be the number of powerful integers in $(n^2,(n+1)^2)$. We record two complementary bounds.
On the lower side we prove a \emph{frequency} (counting) bound, uniform as $V\to\infty$: there are effective
constants $c,C>0$ with $\#\{n\le x:h(n)\ge V\}\ge x\exp(-CV\log(2V)\log\log(3V))$ whenever
$V^2\log(2V)\log\log(3V)\le c\log x$, so in particular $\#\{n\le x:h(n)\ge V\}\ge x^{1-o(1)}$ for
$2\le V\le(\log x)^{1/2-\varepsilon}$. On the upper side we record $h(n)\ll_\varepsilon n^{6/25+\varepsilon}$,
improving the elementary $n^{2/5}$ and the recorded unconditional $O(n/\log n)$. The construction underlying
the lower bound is that of De Koninck--Luca--Shparlinski with the standard primorial modulus and is not
claimed new; the upper bound is an application of the classical integer-points-close-to-a-curve method of
Swinnerton-Dyer and Filaseta--Trifonov, whose only novelty is the explicit worst-case ($\theta=0$) statement.
The new mathematical point is the reduction of the frequency question to simultaneous equidistribution of the
multiquadratic directions $\{1/\sqrt d\}$, made effective by a Liouville bound. The elementary and algebraic
core of the lower bound is formalized in Lean~4/Mathlib. Both sides leave a gap to the conjectured
$(\log n)^{1+o(1)}$, governed in each case by an equidistribution input beyond current technology.
\end{abstract}

\section{Introduction}
A positive integer is \emph{powerful} (squarefull) if every prime dividing it does so to power at least
two. Let $h(n)=\#\{\text{powerful }m: n^2<m<(n+1)^2\}$. (Erd\H{o}s's formulation counts the left endpoint
$n^2$; our open-interval normalization differs by exactly one, as $n^2$ is always powerful, which is
immaterial for the estimates below.) Throughout, $\norm x$ denotes the distance from $x\in\R$ to the nearest
integer, $\norm a_\infty=\max_j|a_j|$, and $e(t)=e^{2\pi it}$. Erd\H{o}s asked for the order of $h(n)$ and,
in the notation of the question, conjectured a constant $c$ with $h(n)\le(\log n)^{c+o(1)}$ for all $n$ and
$h(n)\ge(\log n)^{c-o(1)}$ infinitely often, $c=1$ being the natural guess \cite{Erdos76,Erdos942}.

\emph{Lower bounds, infinitely often.} De Koninck--Luca \cite{DeLu04} proved
$h(n)\gg(\log n/\log\log n)^{1/3}$ for infinitely many $n$; De Koninck--Luca--Shparlinski \cite{DLS} extended
the method to $\kappa$-full integers in $(n^\kappa,(n+1)^\kappa)$, with exponent $1/3+o(1)$; Blomer (recorded
in \cite[Appendix~A]{MSY}) improved the exponent to $1/2$. Applying Dirichlet's theorem directly and taking
the auxiliary modulus to be a primorial sharpens the infinitely-often exponent further to $1-o(1)$, so that
any admissible $c$ above satisfies $c\ge1$; this is a folklore optimization of the same construction (the
primorial choice goes back to Ramanujan), recorded by the author at \cite{Erdos942}, which we do not claim as
new and mention only for context, as it is not used below.

\emph{Distribution, fixed $\ell$.} Shiu \cite{Shiu}, and Xiong--Zaharescu \cite{XZ} for higher powers, proved
that the density $\delta_\ell=\operatorname{dens}\{n:h(n)=\ell\}$ exists for each fixed $\ell$, with
generating function $\sum_\ell\delta_\ell z^\ell=\prod_\lambda(1+(z-1)/\lambda)$ over an explicit set;
Narumi--Tachiya \cite{NT} treat consecutive intervals. These results are not uniform in $\ell$ and give no
information on $\#\{n\le x:h(n)\ge V\}$ when $V=V(x)\to\infty$.

\emph{Upper bounds.} The unconditional worst-case bounds in the literature are short-interval bounds: De
Koninck--Luca--Shparlinski \cite{DLS} and Chan \cite{Chan} give $O(y/\log y)$ powerful numbers in an interval
of length $y$, which at $y\asymp\sqrt x$ (our case, $x=n^2$) is $O(n/\log n)$; an $n^{1/2}$-type bound holds
only under the $abc$ conjecture \cite[Thm.~2]{DLS}. Sawin observed \cite{Erdos942} that the elementary
estimate $h(n)\ll n^{2/5}$ follows from the representation of powerful numbers as $a^2b^3$. The distribution
of squarefull numbers in short intervals of length $x^{1/2+\theta}$ with $\theta>0$ has been studied
intensively in the \emph{asymptotic} regime \cite{Shiu,FT94,HT96,Tri02}, the current record being
$\theta>19/154$ \cite{Tri02}; but the \emph{worst-case} count at the critical length $y\asymp\sqrt x$
($\theta=0$) does not appear to have been recorded.

We supply a frequency lower bound (Theorem~\ref{thm:main}), and record the upper bound that the standard
near-curve technology yields at the critical length (Theorem~\ref{thm:upper}).

\begin{theorem}\label{thm:main}
There are absolute, effectively computable constants $c,C>0$ such that the following holds. Let
$H(V;x)=\#\{n\le x:h(n)\ge V\}$. If $x$ is large and $2\le V$ and
$V^2\log(2V)\log\log(3V)\le c\log x$, then
\[
  H(V;x)\ \ge\ x\exp\!\bigl(-C\,V\log(2V)\log\log(3V)\bigr).
\]
\end{theorem}

\begin{corollary}\label{cor:half}
For every fixed $\varepsilon>0$, $H(V;x)\ge x^{1-o(1)}$ uniformly for $2\le V\le(\log x)^{1/2-\varepsilon}$.
\end{corollary}

\begin{theorem}\label{thm:upper}
For every $\varepsilon>0$, $h(n)\ll_\varepsilon n^{6/25+\varepsilon}$.
\end{theorem}

To the author's knowledge Theorem~\ref{thm:main} is the first counting lower bound for $h(n)$ valid as
$V\to\infty$; the existing distribution results are for fixed $\ell$. Theorem~\ref{thm:upper} is not a new
theorem in method: it is an application of the integer-points-close-to-a-curve technology of Swinnerton-Dyer
\cite{SD} and Filaseta--Trifonov \cite{FT96,Tri98}; we record it because the literature optimized the
asymptotic ($\theta>0$) regime and did not state the worst-case exponent, and because it answers the question
of how far Sawin's elementary $n^{2/5}$ can be pushed by classical means (\S\ref{sec:floor}).

\section{The construction and an exact one-sided window criterion}
Let $D=p_1\cdots p_h$ be the product of the first $h$ primes and let $d_1,\dots,d_\ell$ ($\ell=2^h-1$) be the
squarefree divisors of $D$ exceeding $1$. For $q\ge1$ put $n=Dq$, and for each $d=d_j$ set
\[
  \alpha_d=d^{-1/2},\qquad r_d=\lceil q\alpha_d\rceil,\qquad m_d=d(Dr_d)^2 .
\]

\emph{Powerfulness.} Each $m_d=dD^2r_d^2$ is powerful: a prime dividing $r_d$ occurs through $r_d^2$ to
exponent $\ge2$; a prime dividing $d$ also divides $D$, hence occurs in $dD^2$ to exponent $\ge3$; a prime
dividing $D$ but not $d$ occurs in $D^2$ to exponent $\ge2$. Distinct $d$ give distinct $m_d$, since the
squarefree kernel of $m_d$ is exactly $d$.

\begin{lemma}\label{lem:window}
For $d>1$, $q/\sqrt d$ is irrational, so $\eta_d:=r_d-q/\sqrt d\in(0,1)$ and $\eta_d=1-\{q/\sqrt d\}$. With
$t_d=D\sqrt d\,\eta_d$,
\[
  m_d-n^2=D^2\bigl(dr_d^2-q^2\bigr)=2nt_d+t_d^2,
\]
and hence
\[
  n^2<m_d<(n+1)^2 \iff 2nt_d+t_d^2<2n+1 \iff t_d<1 \iff \Bigl\{\tfrac{q}{\sqrt d}\Bigr\}>1-\tfrac1{D\sqrt d}.
\]
\end{lemma}
\begin{proof}
$m_d-n^2=D^2(q+\sqrt d\,\eta_d)^2-D^2q^2=2D^2q\sqrt d\,\eta_d+D^2d\eta_d^2=2nt_d+t_d^2$. As
$(n+1)^2-n^2=2n+1$ and $t_d\ge0$, $2nt_d+t_d^2<2n+1\iff t_d<1$; and $t_d<1\iff\eta_d<1/(D\sqrt d)\iff
\{q/\sqrt d\}>1-1/(D\sqrt d)$.
\end{proof}

The ceiling $r_d=\lceil q\alpha_d\rceil$ gives a clean one-sided criterion with no tie or side ambiguity.

\section{Reduction to a box-hit count}
Put $\delta_d=1/(D\sqrt d)$, $\alpha=(d^{-1/2})_{d\mid D,d>1}\in[0,1)^\ell$, and
$B=\prod_{d\mid D,d>1}(1-\delta_d,1)\subset[0,1)^\ell$. By Lemma~\ref{lem:window}, if $q\alpha\bmod1\in B$ then
$(n^2,(n+1)^2)$ contains the $\ell$ distinct powerful integers $m_d$, so $h(Dq)\ge\ell$. Hence, for
$Q=\lfloor x/D\rfloor$,
\[
  \#\{n\le x:h(n)\ge\ell\}\ \ge\ \#\{q\le Q: q\alpha\bmod1\in B\}.
\]
Since each prime of $D$ divides exactly half of its $2^h$ divisors, $\prod_{d\mid D,d>1}d=D^{(\ell+1)/2}$, so
\[
  \operatorname{vol}(B)=\prod_{d\mid D,d>1}\frac1{D\sqrt d}=D^{-\ell}\Bigl(\prod_{d}d\Bigr)^{-1/2}=D^{-(5\ell+1)/4}.
\]

\section{A multiquadratic Liouville bound}
\begin{lemma}\label{lem:liouville}
For $a\in\Z^\ell\setminus\{0\}$ with $H=\norm{a}_\infty$, $\norm{a\cdot\alpha}\ge A^{-1}H^{-\ell}$, where
$A=D(3D\ell)^\ell$.
\end{lemma}
\begin{proof}
$\gamma:=D(a\cdot\alpha)=\sum_{d}a_d(D/d)\sqrt d$ is an algebraic integer of
$K=\Q(\sqrt{p_1},\dots,\sqrt{p_h})$, $[K:\Q]=2^h=\ell+1$. Let $m$ be the nearest integer to $a\cdot\alpha$.
Then $\gamma-Dm$ is a nonzero algebraic integer (nonzero because $\{1\}\cup\{\sqrt d:d\mid D,\,d>1\}$ is
$\Q$-linearly independent, so $a\cdot\alpha\notin\Z$), whence $|\Nm_{K/\Q}(\gamma-Dm)|\ge1$. With $S=\sum_d d^{-1/2}\le
\ell/\sqrt2$ and $|m|\le HS+\tfrac12$, every conjugate satisfies
$|\sigma(\gamma-Dm)|\le\sum_d|a_d|(D/d)\sqrt d+D|m|\le HDS+D(HS+\tfrac12)\le3HD\ell=:M$. Separating the
identity embedding from the other $\ell$ gives $1\le|\gamma-Dm|\,M^\ell$, so
$\norm{a\cdot\alpha}=|\gamma-Dm|/D\ge D^{-1}(3D\ell)^{-\ell}H^{-\ell}$.
\end{proof}
The exponent is $\ell$, not $2^\ell$: the field is generated by $h=\log_2(\ell+1)$ square roots, so its degree
is $\ell+1$. The bound is fully effective.

\section{An effective discrepancy estimate}\label{sec:disc}
\begin{proposition}\label{prop:disc}
There is an absolute $C_0>0$ such that for every $Q\ge3$ the star-discrepancy of $(q\alpha)_{q\le Q}$ in
$[0,1)^\ell$ satisfies $D_Q^\ast\le C_0^\ell\,D\ell\,\log(2Q)\,Q^{-1/(\ell+1)}$. Consequently there is an
absolute $C_1>0$ such that if $\log Q\ge C_1\,\ell^2\log(3D)$ then
$\#\{q\le Q:q\alpha\bmod1\in B\}\ge\tfrac12\,Q\operatorname{vol}(B)$.
\end{proposition}
\begin{proof}
By the Erd\H{o}s--Tur\'an--Koksma inequality \cite{DT}, for every integer $R\ge2$,
\[
  D_Q^\ast\le C^\ell\Bigl(\tfrac1R+\sum_{0<\norm a_\infty\le R}\tfrac1{r(a)}\bigl|\tfrac1Q\textstyle\sum_{q\le Q}e(q\,a\cdot\alpha)\bigr|\Bigr),\quad r(a)=\textstyle\prod_j\max(1,|a_j|),
\]
with $C$ absolute. The geometric sum gives $|\sum_{q\le Q}e(q\,a\cdot\alpha)|\le1/(2\norm{a\cdot\alpha})$, and
by Lemma~\ref{lem:liouville} $\norm{a\cdot\alpha}^{-1}\le D(3D\ell)^\ell R^\ell$ for $\norm a_\infty\le R$.
Also $\sum_{\norm a_\infty\le R}r(a)^{-1}\le(1+2\sum_{m\le R}m^{-1})^\ell\le(C\log2R)^\ell$. Hence
$D_Q^\ast\le C^\ell(R^{-1}+D(3D\ell)^\ell R^\ell(\log2R)^\ell/Q)$. Take
$R=\lfloor(Q/[D(3D\ell)^\ell(\log2Q)^\ell])^{1/(\ell+1)}\rfloor$ when this is $\ge2$ (otherwise
$D_Q^\ast\le1$ trivially, absorbed by enlarging $C_0$). Both terms are then
$\ll C^\ell(D(3D\ell)^\ell(\log2Q)^\ell/Q)^{1/(\ell+1)}$, and since
$(D(3D\ell)^\ell)^{1/(\ell+1)}\le3D\ell$, $D_Q^\ast\le C_0^\ell D\ell\log(2Q)Q^{-1/(\ell+1)}$. The box
$B=\prod_j(1-\delta_j,1)$ is anchored at a corner, so by inclusion--exclusion the counting error is
$\le2^\ell QD_Q^\ast$, giving $\#\{q\le Q:q\alpha\in B\}=Q\operatorname{vol}(B)+O(2^\ell QD_Q^\ast)$. This
error is below $\tfrac12Q\operatorname{vol}(B)=\tfrac12QD^{-(5\ell+1)/4}$ once
$2^\ell C_0^\ell D\ell\log(2Q)Q^{-1/(\ell+1)}\le\tfrac12D^{-(5\ell+1)/4}$, i.e.\ (taking logarithms) once
$L/(\ell+1)\ge\tfrac{5\ell+5}4\log D+O(\ell\log\ell)+O(\log L)$, where $L=\log Q$. Since
$L/(\ell+1)-O(\log L)$ is increasing in $L$ once $L\gg\ell$, it suffices to verify this at
$L=C_1\ell^2\log(3D)$, which holds for an absolute $C_1$.
\end{proof}

\section{Proof of Theorem~\ref{thm:main}}
Choose $h$ minimal with $2^h-1\ge V$ and set $\ell=2^h-1$, $D=p_1\cdots p_h$, so $V\le\ell\le2V+1$. By
standard explicit (Chebyshev-type) prime estimates $\log D=\theta(p_h)\ll h\log(2h)\ll\log(2V)\log\log(3V)$,
hence $\ell^2\log(3D)\ll V^2\log(2V)\log\log(3V)$. The hypothesis $V^2\log(2V)\log\log(3V)\le c\log x$ gives,
for $c$ small, $\log(x/D)\ge C_1\ell^2\log(3D)$. Applying Proposition~\ref{prop:disc} with
$Q=\lfloor x/D\rfloor\ge x/(2D)$,
\[
  H(V;x)\ge\#\{n\le x:h(n)\ge\ell\}\ge\tfrac12 Q\operatorname{vol}(B)\ge\tfrac{x}{4}\,D^{-(5\ell+5)/4}.
\]
Since $\ell\log D\ll V\log(2V)\log\log(3V)$, this is $\ge x\exp(-C\,V\log(2V)\log\log(3V))$.
Corollary~\ref{cor:half} follows because $V\log(2V)\log\log(3V)=o(\log x)$ for
$V\le(\log x)^{1/2-\varepsilon}$.\qed

\section{A conditional near-ceiling form}
\begin{theorem}\label{thm:cond}
Let $D=p_1\cdots p_h$, $\ell=2^h-1$, $Q=\lfloor x/D\rfloor$, and assume the simultaneous-equidistribution
hypothesis $D_Q^\ast(\alpha)\le2^{-\ell-1}\operatorname{vol}(B)$. Then $\#\{n\le x:h(n)\ge\ell\}\ge
\tfrac12\lfloor x/D\rfloor D^{-(5\ell+1)/4}$, which for $x\ge2D$ is
\[
  \#\{n\le x:h(n)\ge\ell\}\ \ge\ \tfrac14\,x\,D^{-(5\ell+5)/4}\ =\ x^{\,1-\frac{5\ell+5}{4}\frac{\log D}{\log x}+o(1)}.
\]
This is a positive power of $x$ while $\tfrac{5\ell}{4}\log D<(1-o(1))\log x$ (below the construction's height
ceiling $\ell^\ast(x)\asymp\log x/(\log\log x\,\log\log\log x)$); it is $x^{1-o(1)}$ only for
$\ell\log D=o(\log x)$, and at $\ell=\ell^\ast$ it degenerates to the infinitely-often statement.
\end{theorem}
\begin{proof}
Under the hypothesis the corner-box counting error is $2^\ell QD_Q^\ast\le\tfrac12Q\operatorname{vol}(B)$, so
$\#\{q\le Q:q\alpha\in B\}\ge\tfrac12Q\operatorname{vol}(B)=\tfrac12\lfloor x/D\rfloor D^{-(5\ell+1)/4}$, exactly as
in the second half of Proposition~\ref{prop:disc}; the displayed forms then follow as in \S\ref{sec:disc}--\S6.
\end{proof}

\section{The upper bound}\label{sec:upper}
We deduce Theorem~\ref{thm:upper} from the integer-points-close-to-a-curve method. Every powerful $m$ has a
unique representation $m=a^2b^3$ with $b$ squarefree. The method we use is the following.

\begin{lemma}[Filaseta--Trifonov {\cite[Thm.~4.1]{FT96}}]\label{lem:ft}
Let $k\ge2$ be an integer and $s\in\Q\setminus\{-(k-1),\dots,k-1\}$. Let $f(u)=X/u^s$ with $X$ real
independent of $u$, and suppose $N^s\le X$ and $\delta\le c_k N^{-(k-1)}$ for a sufficiently small
$c_k=c_k(k,s)>0$. Then
\[
 \#\{u\in\Z\cap(N,2N]:\norm{f(u)}<\delta\}\ \ll_{k,s}\ X^{1/(2k+1)}N^{(k-s)/(2k+1)}+\delta\,X^{1/(6k+3)}N^{(6k^2+2k-s-1)/(6k+3)}.
\]
\end{lemma}

\begin{proof}[Proof of Theorem~\ref{thm:upper}]
Write $m=a^2b^3$ with $b$ squarefree. If $n^2<m<(n+1)^2$ then $a^2b^3<(n+1)^2$, so $a<n+1$,
$b<(n+1)^{2/3}$, and $\min(a,b)\le(n+1)^{2/5}$ (else $a,b>(n+1)^{2/5}$ give $a^2b^3>(n+1)^2$). We bound the
contributions of $b\le(n+1)^{2/5}$ and of $a\le(n+1)^{2/5}$ separately; double counting only helps an upper
bound, and the $O(1)$ terms with small parameter $\le2$ are discarded. Dropping the squarefree restriction on
$b$ likewise only enlarges the count.

\emph{The $b$-aspect.} For fixed $b\ge2$, the condition $n^2<a^2b^3<(n+1)^2$ places $a$ in an interval
$(nb^{-3/2},(n+1)b^{-3/2})$ of length $b^{-3/2}<1$; so there is at most one admissible $a$, and its existence
forces $\norm{nb^{-3/2}}<b^{-3/2}$. Thus the $b\le(n+1)^{2/5}$ contribution is
$\le\#\{b\le(n+1)^{2/5}:\norm{nb^{-3/2}}<b^{-3/2}\}$. Cover $b\in(2,(n+1)^{2/5}]$ by dyadic blocks
$b\in(B,2B]$, $B=n^\beta$, $0<\beta\le\tfrac25+o(1)$, and apply Lemma~\ref{lem:ft} with $k=2$, $s=3/2$, $X=n$,
$N=B$, $\delta\asymp B^{-3/2}$. The hypotheses hold: $s=3/2\notin\{-1,0,1\}$; $N^s=B^{3/2}\le n$ since
$B\le n^{2/3}$; and $\delta\asymp B^{-3/2}\le c_2B^{-1}$ for $B\ge c_2^{-2}$ (boundedly many smaller $B$
contribute $O(1)$). With $6k^2+2k-s-1=51/2$, the bound is
\[
  \ll n^{1/5}B^{1/10}+B^{-3/2}n^{1/15}B^{51/30}=n^{1/5+\beta/10}+n^{1/15+\beta/5},
\]
the first term dominating for $\beta\le2/3$. Together with the trivial block bound $\#\{b\in(B,2B]\}\ll B=n^\beta$,
the per-block count is $\ll\min(n^\beta,\,n^{1/5+\beta/10})$, which on $0<\beta\le\tfrac25+o(1)$ is maximized at
$\beta=\tfrac25$, giving $n^{6/25+o(1)}$.

\emph{The $a$-aspect.} For fixed $a$, $b$ lies in $((n/a)^{2/3},((n+1)/a)^{2/3})$, an interval of length
$\asymp a^{-2/3}n^{-1/3}<1$, so at most one admissible $b$, requiring $\norm{(n/a)^{2/3}}\ll a^{-2/3}n^{-1/3}$.
Cover $a\in(2,(n+1)^{2/5}]$ dyadically, $a\in(A,2A]$, $A=n^\alpha$, $0<\alpha\le\tfrac25+o(1)$, and apply
Lemma~\ref{lem:ft} with $k=2$, $s=2/3$, $X=n^{2/3}$, $N=A$, $\delta\asymp A^{-2/3}n^{-1/3}$. Validity:
$N^s=A^{2/3}\le n^{2/3}=X$ since $A\le n$; and $\delta A/c_2\asymp A^{1/3}n^{-1/3}\ll n^{-1/5}$ since
$A\le(n+1)^{2/5}$, so $\delta\le c_2A^{-1}$ for $n$ large (boundedly many small $n$ are absorbed). The leading
term is $X^{1/5}N^{(k-s)/5}=n^{2/15}A^{4/15}=n^{2/15+4\alpha/15}$, maximized at $\alpha=\tfrac25$, giving
$n^{2/15+8/75}=n^{6/25+o(1)}$ (the second term is smaller).

Summing over the $O(\log n)$ dyadic blocks in each aspect costs a factor $n^{o(1)}$. Hence
$h(n)\ll_\varepsilon n^{6/25+\varepsilon}$.
\end{proof}

\section{The floor of the method, and the gap}\label{sec:floor}
The exponent $6/25$ is attained at the symmetric scale $a\asymp b\asymp n^{2/5}$, where $m=a^2b^3$ has both
factors at the critical size. There, in the normalization of Lemma~\ref{lem:ft}, the relevant window has
half-width $\asymp n^{-3/5}$ and the curve has $|f''|\asymp n^{-2/5}$, so $6/25$ coincides with the
on-curve lattice-point bound $N^{3/5}$ of Swinnerton-Dyer \cite{SD} at $N=n^{2/5}$.

\begin{remark}\label{rem:floor}
For the Filaseta--Trifonov/Swinnerton-Dyer route used here, $6/25$ is the best attainable exponent. A smaller
exponent would require either the higher-order ($k\ge3$) cases of Lemma~\ref{lem:ft}, which are inadmissible
at the symmetric scale because the window half-width $\asymp n^{-3/5}$ exceeds the admissibility threshold
$c_kN^{-(k-1)}=c_3n^{-4/5}$ there; or an estimate of Bombieri--Pila on-curve quality
($\ll N^{1/2+\varepsilon}$, which would give $n^{1/5}$) valid for a positive window half-width in the regime
where the two curve scales coincide. Trifonov \cite{Tri98,Tri02} identifies exactly this --- extending
Swinnerton-Dyer's on-curve count to the positive-width case at the symmetric scale $a\asymp b\asymp n^{2/5}$
(both factors of size $\asymp X^{1/5}$ for squarefull numbers near $X=n^2$) --- as the open obstacle, and no
such extension is known. Thus $6/25$ is the natural limit of present technology for the worst-case count,
consistent with Sawin's expectation \cite{Erdos942} that arguments of this kind yield a power of $n$ rather
than a power of $\log n$.
\end{remark}

\begin{remark}
Numerically $\max_{n\le10^7}h(n)=9$, attained at $n=524827$, and the growth of $\max_{n\le N}h(n)$ is
consistent with $(\log N)^{1+o(1)}$. The gap between this conjectured truth and our $n^{6/25}$ on the upper
side, and between it and the height ceiling $\ell^\ast(x)$ on the lower side (Theorem~\ref{thm:cond}), is in
each case governed by an equidistribution input beyond current technology: a positive-width on-curve count on
the upper side, and a lower bound for $\norm{\sum_d a_d\sqrt d}$ uniform in $\ell$ on the lower side. These
appear related to the arithmetic of the conjectured bound $h(n)\le(\log n)^{1+o(1)}$, though that bound could
fail for other reasons.
\end{remark}

\begin{remark}
The lower-bound construction is a quantitative refinement of \cite{DeLu04,DLS} with the standard primorial
modulus; the new point is not the construction but the extraction of a counting lower bound for $\#\{n\le
x:h(n)\ge V\}$ uniform as $V\to\infty$, via an exact reduction to a box-hit count (\S2--\S3), a discrepancy
criterion (Proposition~\ref{prop:disc}, second half), and the quantitative Erd\H{o}s--Tur\'an--Koksma input
made effective by Lemma~\ref{lem:liouville}. The upper bound (\S\ref{sec:upper}) uses only classical
near-curve technology and is recorded for the worst-case length, where the literature is silent.
\end{remark}

\begin{remark}
The elementary core of the lower bound --- powerfulness of the $m_d$, the exact window criterion
Lemma~\ref{lem:window}, distinctness, and the deduction $h(Dq)\ge\ell$ from a box hit --- is formalized in
Lean~4/Mathlib with no \texttt{sorry} on the standard axioms $\{$\texttt{propext}, \texttt{Classical.choice},
\texttt{Quot.sound}$\}$, as is the abstract Liouville norm mechanism. The multiquadratic Liouville inequality
(Lemma~\ref{lem:liouville}) and the quantitative Erd\H{o}s--Tur\'an--Koksma input (Proposition~\ref{prop:disc})
are each recorded as one explicit named axiom, since neither is presently in Mathlib (both are classical,
provable mathematics). Thus the formalization verifies the algebraic and elementary core and the final counting
deduction \emph{modulo} two stated classical inputs. The Lean development is available in the accompanying
repository, built against Mathlib v4.30.0.
\end{remark}

\section*{Use of AI}
Claude (Opus~4.8), Codex (GPT-5.5) and Aristotle were used in this work, including for adversarial checking of
the arguments and for the Lean formalization; the author has checked the mathematics and the cited literature
and takes responsibility for all claims. The Lean kernel checks the elementary and algebraic core of the lower
bound on the standard Mathlib axioms; the formal counting theorem is checked relative to the two explicit named
axioms recorded above.

\begin{thebibliography}{99}
\bibitem{Erdos76} P.~Erd\H{o}s, \emph{Problems and results on number theoretic properties of consecutive integers and related questions}, Proc.\ Fifth Manitoba Conf.\ on Numerical Math.\ (Univ.\ Manitoba, 1975), Congress.\ Numer.\ XVI, Utilitas Math., 1976, pp.~25--44.
\bibitem{Erdos942} T.~F. Bloom (ed.), \emph{Erd\H{o}s Problem \#942}, \url{https://www.erdosproblems.com/942} (with comments by S.~D. Hughes, N.~Sothanaphan, T.~F. Bloom, and W.~Sawin).
\bibitem{DeLu04} J.-M.~De Koninck and F.~Luca, \emph{Sur la proximit\'e des nombres puissants}, Acta Arith. \textbf{114} (2004), no.~2, 149--157.
\bibitem{DLS} J.-M.~De Koninck, F.~Luca and I.~E.~Shparlinski, \emph{Powerful numbers in short intervals}, Bull. Austral. Math. Soc. \textbf{71} (2005), no.~1, 11--16.
\bibitem{MSY} M.~Munsch, I.~E.~Shparlinski and K.~H.~Yau, \emph{Smooth squarefree and squarefull integers in arithmetic progressions}, Mathematika \textbf{66} (2020), no.~1, 56--70. (Appendix~A, due to V.~Blomer, improves the infinitely-often exponent from $1/3$ to $1/2$.)
\bibitem{Shiu} P.~Shiu, \emph{On the number of square-full integers between successive squares}, Mathematika \textbf{27} (1980), no.~2, 171--178.
\bibitem{XZ} M.~Xiong and A.~Zaharescu, \emph{$k$-full integers between successive $k$-th powers}, Indag. Math. (N.S.) \textbf{22} (2011), no.~1--2, 77--86.
\bibitem{NT} S.~Narumi and Y.~Tachiya, \emph{On the number of $k$-full integers between three successive $k$-th powers}, arXiv:2512.07438 (2025).
\bibitem{Chan} T.~H.~Chan, \emph{A note on powerful numbers in short intervals}, Bull. Aust. Math. Soc. \textbf{108} (2023), no.~1, 99--106.
\bibitem{SD} H.~P.~F.~Swinnerton-Dyer, \emph{The number of lattice points on a convex curve}, J. Number Theory \textbf{6} (1974), 128--135.
\bibitem{FT94} M.~Filaseta and O.~Trifonov, \emph{The distribution of squarefull numbers in short intervals}, Acta Arith. \textbf{67} (1994), no.~4, 323--333.
\bibitem{FT96} M.~Filaseta and O.~Trifonov, \emph{The distribution of fractional parts with applications to gap results in number theory}, Proc. London Math. Soc. (3) \textbf{73} (1996), no.~2, 241--278.
\bibitem{HT96} M.~N.~Huxley and O.~Trifonov, \emph{The square-full numbers in an interval}, Math. Proc. Cambridge Philos. Soc. \textbf{119} (1996), no.~2, 201--208.
\bibitem{Tri98} O.~Trifonov, \emph{Integer points close to a smooth curve}, Serdica Math. J. \textbf{24} (1998), no.~3--4, 319--338.
\bibitem{Tri02} O.~Trifonov, \emph{Lattice points close to a smooth curve and squarefull numbers in short intervals}, J. London Math. Soc. (2) \textbf{65} (2002), no.~2, 303--319.
\bibitem{DT} M.~Drmota and R.~F.~Tichy, \emph{Sequences, Discrepancies and Applications}, Lecture Notes in Math. \textbf{1651}, Springer, 1997.
\end{thebibliography}

\bigskip
\noindent\textsc{Scott D. Hughes, Independent Researcher}\\
\noindent\textit{Email address:} \texttt{v@deltagray.com}

\end{document}
